% GENERATED FILE. Edit canonical lessons, then run npm run generate:books.
% canonical-source-hash: fnv1a64:eb8af5ebdd595bc2
% canonical-lessons: SA-C300-bhukampah, SA-C300-kampate, SA-C300-isat, SA-C300-sphutati, SA-C300-gartah

\chapter{When the Ground Shook}
\label{ch:sa-when-the-ground-shook}
\hlchaptermodality{300}

\begin{chapteropening}
\textbf{By the end of this chapter:} \emph{I can tell what happened in an earthquake.}

Complete the last lesson of chapter 300: I can tell what happened in an earthquake.
\end{chapteropening}

\section[\texorpdfstring{\sk{भूकम्पः} (bhūkampaḥ)}{भूकम्पः (bhūkampaḥ)}]{\sk{भूकम्पः} (bhūkampaḥ) — an earthquake}

\label{lesson:SA-C300-bhukampah}



\begin{quote}
Before the new one: say the Sanskrit for keeps away, then the Sanskrit for a misfortune.
\end{quote}

\subsection*{You'll want to know: \sk{भूकम्पः}}
\textbf{\sk{भूकम्पः}} — \emph{bhūkampaḥ} — \textquotedblleft{}an earthquake\textquotedblright{}.

\begin{quote}
\sk{ह्यः} \sk{रात्रौ} \sk{भूकम्पः} \sk{अभवत्।}
\end{quote}

\begin{quote}
\emph{hyaḥ rātrau bhūkampaḥ abhavat.}
\end{quote}

\begin{quote}
There was an earthquake last night.
\end{quote}

A masculine noun.

\begin{cousinweb}
\textbf{\sk{भूकम्पः}} (\emph{bhūkampaḥ}) is built from \emph{bhū}, \textquotedblleft{}the earth\textquotedblright{}, and \emph{kampa}, \textquotedblleft{}trembling\textquotedblright{}; MW gives \emph{bhūkampa} as an earthquake (MW).
\end{cousinweb}

\begin{grammarlens}[title={What you've built}]
One more word for telling what happened.
\end{grammarlens}

\subsection*{Your turn}
\begin{itemize}
\raggedright
  \item \emph{Say it:} \emph{bhūkampaḥ}
  \item \emph{Say it:} \emph{bhūkampaḥ}, once more
  \item \emph{Say it:} \emph{hyaḥ rātrau bhūkampaḥ abhavat.} — There was an earthquake last night.
  \item \emph{Recall:} say the Sanskrit for keeps away, then the Sanskrit for a misfortune, then say \emph{bhūkampaḥ} again
\end{itemize}

\begin{tcolorbox}[breakable,colback=teal!4,colframe=teal!35!black,title={Before you move on}]
What does \sk{भूकम्पः} mean? (\textquotedblleft{}An earthquake\textquotedblright{}.) Say it once more.
\end{tcolorbox}
\section[\texorpdfstring{\sk{कम्पते} (kampate)}{कम्पते (kampate)}]{\sk{कम्पते} (kampate) — shakes; trembles}

\label{lesson:SA-C300-kampate}



\begin{quote}
Before the new one: say the Sanskrit for a misfortune, then the Sanskrit for an earthquake.
\end{quote}

\subsection*{You'll want to know: \sk{कम्पते}}
\textbf{\sk{कम्पते}} — \emph{kampate} — \textquotedblleft{}shakes; trembles\textquotedblright{}.

\begin{quote}
\sk{भूमिः} \sk{कम्पते}, \sk{जनाः} \sk{पलायन्ते।}
\end{quote}

\begin{quote}
\emph{bhūmiḥ kampate, janāḥ palāyante.}
\end{quote}

\begin{quote}
The ground shakes and people run away.
\end{quote}

Past (it …): \sk{भूमिः} \sk{कम्पितवती}, \emph{bhūmiḥ kampitavatī}.

\begin{cousinweb}
\textbf{\sk{कम्पते}} (\emph{kampate}) comes from the root √kamp, \textquotedblleft{}to tremble, shake\textquotedblright{}; MW compares Greek \emph{kamptō} (MW).
\end{cousinweb}

\begin{grammarlens}[title={What you've built}]
One more word for telling what happened.
\end{grammarlens}

\subsection*{Your turn}
\begin{itemize}
\raggedright
  \item \emph{Say it:} \emph{kampate}
  \item \emph{Say it:} \emph{kampate}, once more
  \item \emph{Say it:} \emph{bhūmiḥ kampate, janāḥ palāyante.} — The ground shakes and people run away.
  \item \emph{Recall:} say the Sanskrit for a misfortune, then the Sanskrit for an earthquake, then say \emph{kampate} again
\end{itemize}

\begin{tcolorbox}[breakable,colback=teal!4,colframe=teal!35!black,title={Before you move on}]
What does \sk{कम्पते} mean? (\textquotedblleft{}Shakes; trembles\textquotedblright{}.) Say it once more.
\end{tcolorbox}
\section[\texorpdfstring{\sk{ईषत्} (īṣat)}{ईषत् (īṣat)}]{\sk{ईषत्} (īṣat) — slightly}

\label{lesson:SA-C300-isat}



\begin{quote}
Before the new one: say the Sanskrit for an earthquake, then the Sanskrit for shakes.
\end{quote}

\subsection*{You'll want to know: \sk{ईषत्}}
\textbf{\sk{ईषत्}} — \emph{īṣat} — \textquotedblleft{}slightly\textquotedblright{}.

\begin{quote}
\sk{प्रथमं} \sk{भूमिः} \sk{ईषत्} \sk{कम्पते।}
\end{quote}

\begin{quote}
\emph{prathamaṃ bhūmiḥ īṣat kampate.}
\end{quote}

\begin{quote}
At first the ground shakes slightly.
\end{quote}

\begin{cousinweb}
\textbf{\sk{ईषत्}} (\emph{īṣat}) MW gives \emph{īṣat} as an indeclinable meaning a little, slightly, without a derivation, so its origin is uncertain (MW).
\end{cousinweb}

\begin{grammarlens}[title={What you've built}]
One more word for telling what happened.
\end{grammarlens}

\subsection*{Your turn}
\begin{itemize}
\raggedright
  \item \emph{Say it:} \emph{īṣat}
  \item \emph{Say it:} \emph{īṣat}, once more
  \item \emph{Say it:} \emph{prathamaṃ bhūmiḥ īṣat kampate.} — At first the ground shakes slightly.
  \item \emph{Recall:} say the Sanskrit for an earthquake, then the Sanskrit for shakes, then say \emph{īṣat} again
\end{itemize}

\begin{tcolorbox}[breakable,colback=teal!4,colframe=teal!35!black,title={Before you move on}]
What does \sk{ईषत्} mean? (\textquotedblleft{}Slightly\textquotedblright{}.) Say it once more.
\end{tcolorbox}
\section[\texorpdfstring{\sk{स्फुटति} (sphuṭati)}{स्फुटति (sphuṭati)}]{\sk{स्फुटति} (sphuṭati) — cracks open; bursts}

\label{lesson:SA-C300-sphutati}



\begin{quote}
Before the new one: say the Sanskrit for shakes, then the Sanskrit for slightly.
\end{quote}

\subsection*{You'll want to know: \sk{स्फुटति}}
\textbf{\sk{स्फुटति}} — \emph{sphuṭati} — \textquotedblleft{}cracks open; bursts\textquotedblright{}.

\begin{quote}
\sk{भित्तिः} \sk{सहसा} \sk{स्फुटति।}
\end{quote}

\begin{quote}
\emph{bhittiḥ sahasā sphuṭati.}
\end{quote}

\begin{quote}
The wall suddenly cracks open.
\end{quote}

Past (it …): \sk{भित्तिः} \sk{सहसा} \sk{स्फुटितवती}, \emph{bhittiḥ sahasā sphuṭitavatī}.

\begin{cousinweb}
\textbf{\sk{स्फुटति}} (\emph{sphuṭati}) comes from the root √sphuṭ, \textquotedblleft{}to burst or split open with a sound\textquotedblright{} (MW).
\end{cousinweb}

\begin{grammarlens}[title={What you've built}]
One more word for telling what happened.
\end{grammarlens}

\subsection*{Your turn}
\begin{itemize}
\raggedright
  \item \emph{Say it:} \emph{sphuṭati}
  \item \emph{Say it:} \emph{sphuṭati}, once more
  \item \emph{Say it:} \emph{bhittiḥ sahasā sphuṭati.} — The wall suddenly cracks open.
  \item \emph{Recall:} say the Sanskrit for shakes, then the Sanskrit for slightly, then say \emph{sphuṭati} again
\end{itemize}

\begin{tcolorbox}[breakable,colback=teal!4,colframe=teal!35!black,title={Before you move on}]
What does \sk{स्फुटति} mean? (\textquotedblleft{}Cracks open; bursts\textquotedblright{}.) Say it once more.
\end{tcolorbox}
\section[\texorpdfstring{\sk{गर्तः} (gartaḥ)}{गर्तः (gartaḥ)}]{\sk{गर्तः} (gartaḥ) — a hole in the ground; a pit}

\label{lesson:SA-C300-gartah}



\begin{quote}
Before the new one: say the Sanskrit for slightly, then the Sanskrit for cracks open.
\end{quote}

\subsection*{You'll want to know: \sk{गर्तः}}
\textbf{\sk{गर्तः}} — \emph{gartaḥ} — \textquotedblleft{}a hole in the ground; a pit\textquotedblright{}.

\begin{quote}
\sk{मार्गे} \sk{महान्} \sk{गर्तः} \sk{अस्ति।}
\end{quote}

\begin{quote}
\emph{mārge mahān gartaḥ asti.}
\end{quote}

\begin{quote}
There is a big pit in the road.
\end{quote}

A masculine noun.

\begin{cousinweb}
\textbf{\sk{गर्तः}} (\emph{gartaḥ}) MW gives this \emph{garta} as the same word as \emph{karta}, a hollow, hole or cave, and gives it no root, so its origin is uncertain (MW).
\end{cousinweb}

\begin{grammarlens}[title={What you've built}]
One more word for telling what happened.
\end{grammarlens}

\subsection*{Your turn}
\begin{itemize}
\raggedright
  \item \emph{Say it:} \emph{gartaḥ}
  \item \emph{Say it:} \emph{gartaḥ}, once more
  \item \emph{Say it:} \emph{mārge mahān gartaḥ asti.} — There is a big pit in the road.
  \item \emph{Recall:} say the Sanskrit for slightly, then the Sanskrit for cracks open, then say \emph{gartaḥ} again
\end{itemize}

\begin{tcolorbox}[breakable,colback=teal!4,colframe=teal!35!black,title={Before you move on}]
What does \sk{गर्तः} mean? (\textquotedblleft{}A hole in the ground; a pit\textquotedblright{}.) Say it once more.
\end{tcolorbox}
